\chapter{Multiplicities and the local vacuum}
\label{chap:collision-vacuum}

When collision adds generator degrees, two degree-one states combine
to degree two.  Replacing their product by a degree-one singleton
changes the coefficient.
The free commutative example already detects the change.

Let
\[
 F=k[x_0,x_1,\ldots],\qquad Tx_i=x_{i+1},
\]
where \(T\) acts as a derivation and every \(x_i\) has generator
degree one.  The degree-one subspace is the free \(k[T]\)-module
generated by \(x_0\).  In degree two, the map
\[
 x_i x_j\longmapsto
 \tfrac12(t_1^it_2^j+t_1^jt_2^i)
\]
identifies the state space with \(k[t_1,t_2]^{\mathfrak S_2}\).
Translation acts there by multiplication by \(t_1+t_2\).  Consequently
\begin{equation}
 F^{(2)}\simeq k[T,P],\qquad T=t_1+t_2,\quad P=t_1t_2,
 \label{eq:vacuum-weight-two}
\end{equation}
as \(k[T]\)-modules.  This module has the infinite basis
\(1,P,P^2,\ldots\).  Its rank differs from that of \(F^{(1)}\).
The comparison is a comparison of state modules, not of algebras:
the product of two degree-two states lies in degree four.

Multiplicity labels retain this degree change.  They also distinguish
a position that carries the vacuum, of weight zero, from the empty
word, which has no position.

\section{Additive labels on collision diagrams}

Retain the coefficient hypotheses and cochain conventions of
Chapter~\ref{chap:collision-words}.  For the realizations below,
assume also that the exceptional pullbacks commute with the
direct-sum simplicial totalizations being used.  A multiplicity
word is a pair
\(([n],\mathbf m)\), where \(\mathbf m=(m_1,\ldots,m_n)\) and
\(m_i\in\mathbb N\), including zero.  A monotone surjection
\(\alpha:[n]\twoheadrightarrow[r]\) gives the target label
\[
 (\alpha_*\mathbf m)_j=\sum_{\alpha(i)=j}m_i.
\]
A coefficient system on these labelled words has coherent
equivalences
\[
 d_\alpha^!F_{n,\mathbf m}\simeq F_{r,\alpha_*\mathbf m}.
\]
Its positions still lie in \(X^n\).  Only the coefficient labels
have changed.  The empty word carries its own coefficient at
\(n=0\).

The cut product keeps the labels on their respective blocks:
\[
 (F\star_cG)_{n,\mathbf m}
 =\bigoplus_{p=0}^n j_{p,*}j_p^!
 \bigl(F_{p,(m_1,\ldots,m_p)}\boxtimes
       G_{n-p,(m_{p+1},\ldots,m_n)}\bigr).
\]
The common-block proof of associativity still applies.  Positivity
still means vanishing at the empty word.  Thus the number of
positive factors in arity \(n\) remains at most \(n\), even when
some nonempty positions have weight zero.

To apply these labels to residue coefficients, fix a nonnegatively
graded state algebra \(A=\bigoplus_{d\geq0}A^{(d)}\).  Assume that
translation preserves degree and that every residue operation adds
state degrees.  In particular, this is an assumption on the chosen
generator grading, not on conformal weight in an arbitrary vertex
algebra.

Suppose a full residue complex has coefficients
\begin{equation}
 C_n(A)=\bigoplus_{\pi,\mathbf d}V_{\pi,\mathbf d},
 \qquad \mathbf d:\pi\longrightarrow\mathbb N.
 \label{eq:vacuum-residue-coefficients}
\end{equation}
Here \(\pi\) is a partition of \([n]\), and its blocks are ordered
by their largest labels.  The coefficient \(V_{\pi,\mathbf d}\)
contains the states \(A^{(d_B)}\) at its blocks, with shift
\([|\pi|]\).  It includes the complete differential-module
coefficient along the partition diagonal, including normal
derivatives and localization where distinct blocks are separate.
It is not a tensor product of the underlying state spaces alone.

The residue differential merges two blocks and adds their degrees.
Assume its square is zero and that its cut and exceptional-collision
maps are specified coherently.  In particular, under a monotone
collision, the localized partition terms killed by pullback vanish.
The remaining terms identify with the quotient partitions and their
summed block degrees.  These are hypotheses on the full residue
construction in \eqref{eq:vacuum-residue-coefficients}.

For a label \(\mathbf m\), define
\begin{equation}
 C_{n,\mathbf m}(A)
 =\bigoplus_\pi V_{\pi,\,(\sum_{i\in B}m_i)_{B\in\pi}}.
 \label{eq:vacuum-multiplicity-coefficient}
\end{equation}

\begin{proposition}[Multiplicity and collision]
\label{prop:vacuum-multiplicity}
The coefficients \eqref{eq:vacuum-multiplicity-coefficient} inherit
the residue differential and form a coefficient system on additive
multiplicity words.  If the full residue coalgebra is strong, its
labelled version is strong.  These statements preserve the specified
normal derivatives and collision maps.
\end{proposition}
\begin{proof}
Merging blocks \(B\) and \(B'\) replaces their degrees by
\(\sum_{i\in B\cup B'}m_i\).  Hence the differential stays in the
labelled complex.  For a surviving block \(D\) of a quotient
partition, its degree is
\[
 \sum_{j\in D}(\alpha_*\mathbf m)_j
   =\sum_{i\in\alpha^{-1}(D)}m_i.
\]
This is exactly the degree in its source coefficient.  The given
exceptional-collision comparison therefore restricts to the labelled
summands.  Composition of collisions adds the same integers in
either order.  On a cut open, surviving partition blocks lie on
one side or the other.  Their degree assignments and residue
differentials split between the two sides, so the strong cut maps
restrict as well.  All operations on the differential-module
coefficients remain the original ones.
\end{proof}

The label \((1,1)\) collides to \((2)\), not to \((1)\).
Equation~\eqref{eq:vacuum-weight-two} shows why these targets cannot
be interchanged even in the free commutative example.

\section{Reassembling the full coefficient}

Summing labelled complexes over all \(\mathbf m\) duplicates merged
states.  At two positions, a merged state of weight two appears in
the labels \((2,0)\), \((1,1)\) and \((0,2)\).  The full residue
coefficient contains that state once.  A simplicial resolution
identifies these copies while preserving their collision maps.

For a fixed summand of \eqref{eq:vacuum-residue-coefficients}, put
\[
 M(\pi,\mathbf d)=
 \{\mathbf m\in\mathbb N^n:
           \sum_{i\in B}m_i=d_B\text{ for every }B\in\pi\}.
\]
This is a nonempty finite set.  It has
\(\prod_{B\in\pi}\binom{d_B+|B|-1}{|B|-1}\) elements: each block
requires a weak composition of its weight into \(|B|\) parts.
Give it the groupoid structure with exactly one morphism between
any two objects.  Its nerve in degree \(q\) is
\(M(\pi,\mathbf d)^{q+1}\).

Replace each coefficient by
\begin{equation}
 R_qV_{\pi,\mathbf d}=
 \bigoplus_{(\mathbf m^0,\ldots,\mathbf m^q)
                  \in M(\pi,\mathbf d)^{q+1}}V_{\pi,\mathbf d}.
 \label{eq:vacuum-assignment-resolution}
\end{equation}
Faces delete an assignment and degeneracies repeat one.  The
augmentation sends every copy to \(V_{\pi,\mathbf d}\) by the
identity.  A residue merger keeps each assignment tuple and adds
the weights of its two blocks.  Every source assignment satisfies
the merged weight condition.  Thus the original residue map
extends to \eqref{eq:vacuum-assignment-resolution} and commutes
with all faces and degeneracies.

In degree zero this construction is
\[
 R_0C_n(A)=\bigoplus_{\mathbf m\in\mathbb N^n}C_{n,\mathbf m}(A).
\]
The higher degrees impose the identifications absent from this
direct sum.  We realize the simplicial complex by direct-sum
totalization.  A coefficient of internal cochain degree \(p\) in
simplicial degree \(q\) has total degree \(p-q\), and its total
differential is \(d_{\mathrm{coef}}+(-1)^p\partial\), where
\(\partial\) is the alternating sum of faces.

\begin{theorem}[Reassembly of multiplicities]
\label{thm:vacuum-reassembly}
The augmentation induces an equivalence
\[
 |R_\bullet C(A)|\longrightarrow C(A).
\]
It preserves the given cut coproducts and exceptional-collision
maps.  The realized coalgebra is strong when \(C(A)\) is strong.
A specified local-vacuum marking through the weight-zero
subcomplex lifts unchanged, and the empty coefficient is preserved.
These assertions use the direct-sum coefficient category above.
\end{theorem}
\begin{proof}
For each \((\pi,\mathbf d)\), choose an assignment \(\mathbf m_*\)
that places \(d_B\) at the largest label of \(B\) and zero at its
other labels.  Prepending \(\mathbf m_*\) defines an extra
degeneracy of the augmented nerve.  If \(\partial\) is its chain
boundary and \(h\) this extra degeneracy, cancellation of the
alternating faces gives
\[
 \partial h+h\partial=1
\]
on the augmented chain complex.  On the unaugmented complex the
right side is \(1-i\varepsilon\), where \(i\) selects
\(\mathbf m_*\) and \(\varepsilon\) is the augmentation.
After tensoring with a coefficient complex, the total homotopy
on internal degree \(p\) is \((-1)^ph\).

The chosen assignments need not commute with residue mergers.
Filter instead by the number of partition blocks.  This
filtration is finite in arity \(n\), and a merger lowers it.
On its associated graded, the preceding contraction proves the
augmentation is an equivalence.  Induction on the finite
filtration proves the same assertion before passing to the
associated graded.  Equivalently, correcting the contraction by
the merger differential takes at most \(n-1\) successive
corrections.  The unbounded simplicial direction therefore causes
no infinite sum on an element.

On a cut open, a surviving partition separates into partitions of
the two blocks.  Its assignment set is the product of their
assignment sets.  This identifies the corresponding simplicial
coefficient with the degreewise external product.  Geometric
realization preserves that product because tensor products
preserve colimits separately and simplicial colimits are sifted.
The resulting cut map is an equivalence whenever the original
cut map is an equivalence.  The common assignment tuples also
give the iterated-cut coherences.

Under a monotone collision, use the prescribed exceptional
comparison on the surviving partition coefficient and send each
assignment to \(\alpha_*\mathbf m\).  This is a functor between
nonempty groupoids with exactly one morphism between each pair
of objects.  It is therefore an equivalence of groupoids.  Its
realization supplies the collision equivalence.  Associativity
of addition gives compatibility with composite collisions.
The degreewise diagrams \(R_qC\) need not themselves have
exceptional pullback equivalences.  Those equivalences hold
after realization.

At weight zero the only assignment is the zero assignment, so
the simplicial resolution is constant on the weight-zero
subcomplex.  The prescribed vacuum marking therefore lifts
through it.  At the empty word the only assignment is the empty
tuple.  Its constant resolution preserves the empty coefficient.
\end{proof}

The positive-multiplicity sectors omit positions carrying the
vacuum.  The all-ones sector also omits the simplices that identify
different assignments to a merged block.  Both omissions change
the full coefficient.  If a completion in generator weight is
introduced, its interchange with realization requires a separate
limit argument.

\section{Derived specialization of the vacuum generator}

A central field that represents the vacuum after specialization
must satisfy every translated vacuum relation.  Let \(A^\circ\)
be a vertex algebra over \(k\) with a central regular differential
subalgebra
\[
 E=k[e_0,e_1,\ldots],\qquad Te_i=e_{i+1}.
\]
Regular means that these fields have no singular products with
any state.  Their normally ordered products give the displayed
commutative algebra.  Assume \(A^\circ\) is flat over \(E\) for
this central action.  The specialization is
\[
 A=A^\circ/(e_0-1,e_1,e_2,\ldots).
\]
Translation forces the relations after the first one because
\(T1=0\).

Adjoin odd central regular generators \(\eta_i\) of cochain
degree \(-1\).  They form an exterior algebra.  Set
\begin{equation}
 Q(A^\circ)=A^\circ\otimes_k\Lambda(\eta_0,\eta_1,\ldots),
 \quad T\eta_i=\eta_{i+1},\quad
 d\eta_i=e_i-\delta_{i0},\quad dA^\circ=0.
 \label{eq:vacuum-derived-specialization}
\end{equation}
All expressions are finite polynomials in these generators.

\begin{theorem}[Translation-compatible specialization]
\label{thm:vacuum-specialization}
Formula~\eqref{eq:vacuum-derived-specialization} defines a
differential graded vertex algebra.  The map
\[
 Q(A^\circ)\longrightarrow A,\qquad \eta_i\longmapsto0,
\]
is a quasi-isomorphism preserving the unit and translation.
Without the flatness assumption, the same complex computes
\(A^\circ\otimes_E^L k\), and its higher Tor classes remain.
\end{theorem}
\begin{proof}
The images \(e_i-\delta_{i0}\) are central regular fields.  Thus
the displayed differential extends as a degree-one derivation
of every vertex product.  It squares to zero on generators and
commutes with translation because
\[
 dT\eta_i=e_{i+1}=Td\eta_i.
\]
These identities extend to every state by the derivation rule.

Put \(x_i=e_i-\delta_{i0}\).  The complex
\[
 K_E=(E\otimes\Lambda(\eta_i),
           d=\textstyle\sum_i x_i\partial_{\eta_i})
\]
is a free resolution of the augmentation \(E\to k\) that sends
all \(x_i\) to zero.  To verify exactness, let \(N\) be total
polynomial degree, counting each \(x_i\) and each \(\eta_i\)
once.  For a monomial with \(N>0\), define
\begin{equation}
 h=\frac1N\sum_i\eta_i\partial_{x_i},\qquad h(1)=0.
 \label{eq:vacuum-koszul-contraction}
\end{equation}
The graded anticommutator of
\(\sum_i x_i\partial_{\eta_i}\) and
\(\sum_i\eta_i\partial_{x_i}\) is the total Euler operator.
Hence \(dh+hd=1-i\varepsilon\), where \(i:k\to K_E\) includes
constants and \(\varepsilon\) evaluates all generators at zero.
Each sum is finite on a polynomial.  Characteristic zero makes
division by \(N\) valid.

Translation acts by
\[
 T=\sum_i x_{i+1}\partial_{x_i}
       +\sum_i\eta_{i+1}\partial_{\eta_i}.
\]
It preserves \(N\).  In its commutator with the numerator of
\eqref{eq:vacuum-koszul-contraction}, the terms
\(\eta_{i+1}\partial_{x_i}\) cancel after reindexing.
Thus \(Th=hT\).  This is a translation-compatible contraction
of the scalar complex.

The homotopy \(h\) is not \(E\)-linear.  Exactness after tensoring
with \(A^\circ\) follows instead from flatness of \(A^\circ\)
over \(E\).  It gives
\(H^0Q(A^\circ)=A\) and zero cohomology in other degrees.
In general, tensoring the free resolution \(K_E\) computes the
derived tensor product and its Tor groups.  The vertex map to
the ordinary specialization preserves the unit and translation
by its defining formulas.
\end{proof}

The relation \(e_0=1\) generally destroys generator grading.
Multiplicity reassembly applies to the graded object before
specialization.  The differential graded specialization then
acts on its full state algebra, without assuming that the old
weight decomposition survives.

Suppose the full residue functor on differential graded state
algebras uses derived tensor products, localization and diagonal
direct images on each partition coefficient.  Suppose also that
its differential and vacuum maps are functorial for vertex maps.
The quasi-isomorphism of
Theorem~\ref{thm:vacuum-specialization} then induces an equivalence
\[
 C(Q(A^\circ))\longrightarrow C(A)
\]
of vacuum-marked collision coalgebras.  Indeed, it is an
equivalence on each partition coefficient, and the finite
filtration by partition blocks proves the assertion with the
residue differential included.  Theorem~\ref{thm:words-finite-duality}
gives an equivalence of their cut cobars.  This argument does
not choose a module-theoretic splitting of any vacuum extension.

\section{Extending a comparison across the collision support}

A comparison defined on the cut opens must extend across their
closed complement.  A prescribed vacuum map adds a relative
condition to that extension problem.

Let \(j:O\hookrightarrow Y\) be an open embedding with closed
complement \(Z\).  Work in a stable coefficient category with
the localization triangle
\[
 S=R\Gamma_ZQ\longrightarrow Q\longrightarrow j_*j^!Q.
\]
Fix a map \(u:V\to P\).  Write \(R\operatorname{Hom}\) for the
derived cochain complex of maps, and define
\begin{align*}
 M&=R\operatorname{Hom}(P,Q),\\
 B&=R\operatorname{Hom}(V,Q)
       \mathop{\times}^{h}_{R\operatorname{Hom}(j^!V,j^!Q)}
          R\operatorname{Hom}(j^!P,j^!Q),\\
 L&=R\operatorname{Hom}(\operatorname{cofib}(u),S).
\end{align*}
The superscript \(h\) denotes a homotopy pullback.  A point of
\(B\) contains a map on \(V\), a map on \(O\), and a specified
homotopy between their restrictions to \(j^!V\).

\Needspace{9\baselineskip}
\begin{theorem}[Relative supported extension]
\label{thm:vacuum-relative-extension}
There is a natural fibre sequence
\[
 L\longrightarrow M\longrightarrow B\longrightarrow L[1].
\]
For compatible degree-zero data \(b\), the obstruction to a
lift with the prescribed homotopy is
\(\partial[b]\in H^1L\).  It vanishes exactly when the space of
such lifts is nonempty.  After choosing one lift, the connected
components of that space form a torsor for \(H^0L\), and its
higher homotopy groups are \(H^{-i}L\), \(i>0\).
\end{theorem}
\begin{proof}
Form the square of mapping complexes obtained by restricting
\(P\to Q\) along \(u\) and along \(j\).  Localization and
adjunction identify its two vertical fibres with
\(R\operatorname{Hom}(P,S)\) and
\(R\operatorname{Hom}(V,S)\).  Its total fibre is therefore
\(R\operatorname{Hom}(\operatorname{cofib}(u),S)=L\).
Taking the other order of fibres identifies the same complex
with the fibre of \(M\to B\).  The long exact cohomology
sequence gives the obstruction.  A lift identifies its fibre
with the additive mapping space of \(L\).  The cochain convention
then gives the stated homotopy groups.
\end{proof}

For an explicit lift, replace \(M\to B\) by a degreewise
surjective mapping-path model with kernel \(L\).  Choose a
graded lift \(m_0\) of a cocycle representing \(b\).  Its
differential \(dm_0\) is a closed element of \(L^1\).  If
\(h\) contracts \(L\), the corrected lift is
\[
 m=m_0-h(dm_0).
\]
It is closed and has the same prescribed data.  Vanishing of
\(H^1L\) supplies existence.  Contractibility of the space of
lifts also requires \(H^iL=0\) for \(i\leq0\).

For strong coalgebras \(F,C\) with vacuum \(J\to F\), suppose
a transverse tensor product \(\otimes^!\) is defined with unit
\(J\).  At configuration size \(n\), take
\[
 P_n=F_n\otimes^!C_n,\qquad V_n=C_n,
 \qquad u_n=({J_n\to F_n})\otimes^!1,
 \qquad Q_n=C_n.
\]
If an action has already been specified on
\(O_n=\bigcup_{p=1}^{n-1}U_p^{(n)}\), compatibly with the vacuum,
its relative extension complex is
\begin{equation}
 L_n=R\operatorname{Hom}
 \left(\operatorname{cofib}(C_n\xrightarrow{u_n}
        F_n\otimes^!C_n),\ R\Gamma_{Z_n}C_n\right).
 \label{eq:vacuum-action-obstruction}
\end{equation}
This complex controls the specified open-and-vacuum extension.
To obtain a coalgebra action on the full collision diagram, one
must also impose its prescribed exceptional restrictions and
all iterated-cut homotopies.  Their relative complex is the
homotopy fibre of the restriction map to that matching diagram.
It cannot be replaced by a product of independent equations.

\section{Higher comparison equations}

A linear comparison must be a map at every configuration size.
An isomorphism of singleton state spaces provides only one of
these components.  Input arity gives a second index: an
\(A_\infty\) Taylor component with \(r\) inputs is a map of
whole coefficient systems, not a component at configuration
size \(r\).

Let \(C,D\) be augmented uncurved \(A_\infty\) algebras in the coefficient
category.  Their structures are degree-one square-zero
coderivations \(b^C,b^D\) on the corresponding tensor coalgebras
of suspended reduced coefficients.  Write \(b_r^C,b_r^D\) for
their components with \(r\) inputs.  Uncurved means \(b_0^C=b_0^D=0\).
A degree-zero coaugmentation-preserving coalgebra map \(F\) is
determined by its components
\(f_r\) with \(r\) inputs.  Assume \(f_1\) is a chain map.
The equation \(b^D F=Fb^C\) in
input arity \(r\geq2\) reads
\begin{equation}
\begin{split}
 \partial f_r={}&
 \sum_{a+s+t=r,\ s\geq2}
 f_{a+1+t}(1^{\otimes a}\otimes b_s^C\otimes1^{\otimes t})\\
 &-\sum_{\substack{k\geq2,\ i_1+\cdots+i_k=r\\i_j\geq1}}
 b_k^D(f_{i_1}\otimes\cdots\otimes f_{i_k})=:O_r.
\end{split}
\label{eq:vacuum-higher-comparison}
\end{equation}
Here \(a,t\geq0\) in the first sum.
The unary mapping differential \(\partial\) is the graded
commutator with the source and target unary differentials.  Every
tensor composition uses the Koszul signs.  Each comparison
component on the right has smaller input arity.  If the lower
equations hold, then \(\partial O_r=0\): expand the arity-\(r\)
part of
\[
 b^D(b^DF-Fb^C)+(b^DF-Fb^C)b^C=0.
\]
Thus \([O_r]\) lies in degree-one cohomology of the appropriate
relative operation complex, including its collision conditions.
A contraction \(h_r\) supplies \(f_r=h_rO_r\) when the
prescribed boundary component is zero.  For a nonzero prescribed
component, first choose a lift and correct its relative defect
as in Theorem~\ref{thm:vacuum-relative-extension}.

By contrast, an unknown structure operation \(b_r\) has degree
one, and the known defect in \(b^2\) has degree two.  Constructing
that operation requires a degree-two obstruction to vanish.
The two extension problems use different cohomology groups.

On positive word coefficients, input arities greater than
configuration size vanish.  The recursion is therefore finite
at each configuration, provided all the relative contractions
and coefficient-domain compatibilities have been supplied.
A nonzero arity-zero operation removes this triangular form:
the arity-\(r\) equation can then contain \(f_{r+1}\).
Such a comparison needs a separate square-zero realization or
a complete positive filtration controlling these insertions.

A surviving obstruction class refutes the comparison whose
domain and prescribed data define that class.  To exclude all
field realizations or all pointed actions, one must prove that
every permitted realization preserves the obstructed datum.
